\documentclass[11pt]{article}
\usepackage[margin=1in]{geometry}
\usepackage{booktabs}
\usepackage{graphicx}
\usepackage{tabularx}
\usepackage{float}
\usepackage{hyperref}
\usepackage{enumitem}
\usepackage{longtable}
\usepackage{array}
\hypersetup{colorlinks=true, linkcolor=blue, urlcolor=blue}
\setlength{\parindent}{0pt}
\setlength{\parskip}{0.6em}

\title{Cross-Dataset Rollout Statistics Report for Qwen/Qwen3-1.7B}
\author{Murphy}
\date{2026-03-16 14:19 UTC}

\begin{document}
\maketitle

\section*{Scope}
This report consolidates the full math $\rightarrow$ GPQA $\rightarrow$ MMLU-Pro $\rightarrow$ LiveCodeBench rollout-stat sweep for \texttt{Qwen/Qwen3-1.7B}. The collector uses the n-gram loop detector with \texttt{n=30} and \texttt{k=20} and records both prompt-level counts and rollout-level event rates, including overlap statistics between looping, max-model-length termination, and correctness.

The math block was evaluated as two separate datasets under the same freeform prompt contract: \texttt{MATH-500} (500 samples) and \texttt{AIME 2024/2025} (60 samples). The downstream multiple-choice and code-generation blocks were evaluated on \texttt{GPQA Diamond}, \texttt{MMLU-Pro}, and \texttt{LiveCodeBench release\_v6} in the same rollout pipeline.

These rows should be read as a common-policy rollout-telemetry bundle under shared low-temperature multi-sample decoding, not as benchmark-optimized single-generation accuracy. The separate GPQA benchmark-style calibration (\texttt{temperature=0.6}, \texttt{num\_generations=1}) answers a different question and should be compared to Qwen's published reference instead of to this cross-dataset table.

\section*{Datasets and Prompt Formats}
\paragraph{MATH-500}
\textbf{Task kind.} \texttt{math\_freeform}
\textbf{Dataset.} 500 free-form math problems from the HuggingFaceH4 MATH-500 test split.
\textbf{Chat / prompt format.} Tokenizer chat template with one user turn: the raw problem text followed by 'You must put your final answer within \textbackslash{}\textbackslash{}boxed\{\}.'

\paragraph{AIME 2024/2025}
\textbf{Task kind.} \texttt{math\_freeform}
\textbf{Dataset.} A local 60-row JSONL made by concatenating the AIME 2024 and 2025 test splits into question/answer records.
\textbf{Chat / prompt format.} The same math-freeform Qwen chat template as MATH-500: one user turn with the problem plus a final-answer-in-box instruction.

\paragraph{GPQA Diamond}
\textbf{Task kind.} \texttt{multiple\_choice\_gpqa}
\textbf{Dataset.} A local staged GPQA Diamond CSV with 198 graduate-level science questions and shuffled four-way answer options.
\textbf{Chat / prompt format.} Tokenizer chat template with one user turn containing the question, a shuffled A-D answer block, and the instruction that the final non-empty line must be a JSON object like '\{"answer": "X"\}'.

\paragraph{MMLU-Pro}
\textbf{Task kind.} \texttt{multiple\_choice\_mmlupro}
\textbf{Dataset.} The TIGER-Lab MMLU-Pro test split, capped to 800 examples for this rollout pass. Requested cap: at most 800 samples.
\textbf{Chat / prompt format.} Tokenizer chat template with one user turn containing the question, an A-J answer list, and the instruction that the final non-empty line must be a JSON object like '\{"answer": "X"\}'.

\paragraph{LiveCodeBench release\_v6}
\textbf{Task kind.} \texttt{livecodebench\_codegen}
\textbf{Dataset.} A question-id-sorted slice of the LiveCodeBench release\_v6 code-generation benchmark, capped to 800 problems for this rollout pass. Requested cap: at most 800 samples.
\textbf{Chat / prompt format.} LiveCodeBench's format\_prompt\_generation pipeline with LM style CodeQwenInstruct, producing raw string prompts rather than a tokenizer chat-template wrapper.

\section*{Generation Configuration}
The common collector configuration across all runs is:
\begin{itemize}[leftmargin=1.5em]
\item Model: \texttt{Qwen/Qwen3-1.7B}
\item Seed: \texttt{0}
\item Loop detector: \texttt{n=30, k=20}
\item Common settings shared by every dataset JSON: \texttt{temperature=0.2}, \texttt{num\_generations=10}, \texttt{max\_tokens=81920}, \texttt{max\_model\_len=40960}, \texttt{dtype=bfloat16}, and \texttt{trust\_remote\_code=True}
\end{itemize}

Per-dataset runtime settings are:
\begin{tabular}{lrrrrr}
\toprule
Dataset & Sample cap & TP & DP & Max seqs & Batch toks \\
\midrule
MATH-500 & full & 1 & 8 & 10 & 4096 \\
AIME 2024/2025 & full & 1 & 8 & 10 & 4096 \\
GPQA Diamond & full & 1 & 4 & 10 & 4096 \\
MMLU-Pro & 800 & 1 & 4 & 10 & 4096 \\
LiveCodeBench release\_v6 & 800 & 1 & 7 & 10 & 4096 \\
\bottomrule
\end{tabular}

\section*{Statistics Logged}
The refreshed collector contract tracks the following metrics across the bundle (13 named metrics in the union):
\begin{itemize}[leftmargin=1.5em]
\item \texttt{avg\_correct\_generation\_length}: Average generated length restricted to correct rollouts.
\item \texttt{avg\_first\_loop\_prefix\_length}: Average prefix length before the first detected loop.
\item \texttt{avg\_generation\_length}: Average number of generated tokens per rollout.
\item \texttt{avg\_loop\_generation\_length}: Average generated length restricted to looped rollouts.
\item \texttt{avg\_wrong\_generation\_length}: Average generated length restricted to wrong rollouts.
\item \texttt{generation\_length\_variance}: Variance of generated token counts across generated rollouts.
\item \texttt{loop\_fraction}: Fraction of generated rollouts flagged by the n-gram loop detector.
\item \texttt{loop\_max\_length\_hit\_fraction}: Among looped rollouts, fraction whose prompt-plus-generation length also hit max\_model\_len.
\item \texttt{loop\_success\_fraction}: Among looped rollouts, fraction that were still correct.
\item \texttt{max\_length\_hit\_fraction}: Fraction of generated rollouts whose prompt-plus-generation length hit max\_model\_len.
\item \texttt{max\_length\_hit\_loop\_fraction}: Among max-length-hit rollouts, fraction that were also looped.
\item \texttt{max\_length\_hit\_success\_fraction}: Among max-length-hit rollouts, fraction that were correct.
\item \texttt{success\_fraction}: Fraction of graded rollouts with the correct answer.
\end{itemize}

The JSON payloads also retain the raw event counts for prompts, graded/generated rollouts, looped rollouts, max-length hits, prompt-too-long exclusions, and their pairwise intersections. For \texttt{LiveCodeBench}, the collector also stores the native benchmark metrics separately from rollout-level success.

\section*{Results Summary}
\begin{tabular}{lrrrrr}
\toprule
Dataset & Prompts & Rollouts & Rollout success & Looped & Max length \\
\midrule
MATH-500 & 500 & 5000 & 72.6\% & 2.9\% & 1.5\% \\
AIME 2024/2025 & 60 & 600 & 38.5\% & 15.0\% & 12.7\% \\
GPQA Diamond & 198 & 1980 & 34.5\% & 16.4\% & 6.9\% \\
MMLU-Pro & 800 & 8000 & 65.2\% & 4.6\% & 0.9\% \\
LiveCodeBench release\_v6 & 800 & 8000 & 58.1\% & 15.0\% & 10.8\% \\
\bottomrule
\end{tabular}

\begin{itemize}[leftmargin=1.5em]
\item MATH-500: rollout success 3628/5000 (72.6\%) graded rollouts; looped 147/5000 (2.9\%) generated rollouts; max-length hits 73/5000 (1.5\%) generated rollouts; looped and max-length 71/147 (48.3\%); max-length and looped 71/73 (97.3\%); looped and correct 47/147 (32.0\%); prompt-too-long 0/500 (0.0\%) prompts.
\item AIME 2024/2025: rollout success 231/600 (38.5\%) graded rollouts; looped 90/600 (15.0\%) generated rollouts; max-length hits 76/600 (12.7\%) generated rollouts; looped and max-length 69/90 (76.7\%); max-length and looped 69/76 (90.8\%); looped and correct 2/90 (2.2\%); prompt-too-long 0/60 (0.0\%) prompts.
\item GPQA Diamond: rollout success 683/1980 (34.5\%) graded rollouts; looped 325/1980 (16.4\%) generated rollouts; max-length hits 137/1980 (6.9\%) generated rollouts; looped and max-length 137/325 (42.2\%); max-length and looped 137/137 (100.0\%); looped and correct 62/325 (19.1\%); prompt-too-long 0/198 (0.0\%) prompts.
\item MMLU-Pro: rollout success 5216/8000 (65.2\%) graded rollouts; looped 365/8000 (4.6\%) generated rollouts; max-length hits 76/8000 (0.9\%) generated rollouts; looped and max-length 76/365 (20.8\%); max-length and looped 76/76 (100.0\%); looped and correct 73/365 (20.0\%); prompt-too-long 0/800 (0.0\%) prompts.
\item LiveCodeBench release\_v6: rollout success 4647/8000 (58.1\%) graded rollouts; looped 1198/8000 (15.0\%) generated rollouts; max-length hits 865/8000 (10.8\%) generated rollouts; looped and max-length 864/1198 (72.1\%); max-length and looped 864/865 (99.9\%); looped and correct 68/1198 (5.7\%); prompt-too-long 0/800 (0.0\%) prompts. Native LiveCodeBench metrics: pass@1 58.1\textbackslash{}\%; pass@5 71.1\textbackslash{}\%; pass@10 73.8\textbackslash{}\%.
\end{itemize}

\subsection*{Headline observations}
The largest raw loop rate appears on GPQA Diamond at 16.4\%, while MATH-500 is the most stable dataset in this bundle at 2.9\%.
The longest looped generations also appear on AIME 2024/2025: 36756.9 tokens on average. The first detected loop appears after 13824.5 tokens on average.
Prompt-plus-generation max-length termination is almost synonymous with looping on the hardest long-form settings: GPQA Diamond reaches 100.0\% for the max-length-hit loop conditional.
The most forgiving loop regime is MATH-500, where 32.0\% of looped rollouts still end in a correct answer.


\subsection*{Recovery caveat}
\begin{itemize}[leftmargin=1.5em]
\item LiveCodeBench release\_v6: the recovered checkpoint contains the final correctness, loop, max-length, and native pass@k metrics, but the original collector crashed after grading and before it wrote an exact \texttt{avg\_first\_loop\_prefix\_length}. That field is intentionally left missing rather than repaired from a non-exact replay.
\end{itemize}


\begin{figure}[H]
\centering
\includegraphics[width=\textwidth]{figures/cross_dataset_rates.png}
\caption{Rollout-level success, loop rate, and max-length-hit rate for each evaluated dataset.}
\end{figure}

\begin{figure}[H]
\centering
\includegraphics[width=\textwidth]{figures/cross_dataset_overlap.png}
\caption{Overlap statistics: among looped rollouts, how many hit max model length; among max-length-hit rollouts, how many looped; and among looped rollouts, how many still answered correctly.}
\end{figure}

\begin{figure}[H]
\centering
\includegraphics[width=\textwidth]{figures/cross_dataset_lengths.png}
\caption{Average generation lengths, average looped-generation lengths, and average first-loop-prefix lengths in thousands of tokens.}
\end{figure}

\section*{Interpretation}
Across the math, science, broad-knowledge, and coding settings, the same pattern repeats: long generation tails are where looping lives. The strongest evidence is in the overlap panel: once a rollout's prompt-plus-generation length reaches \texttt{max\_model\_len}, it is often also a detected loop. The converse is weaker but still substantial: in several datasets, a large fraction of looped rollouts also terminate at that full-context ceiling.

The length panel shows that looped generations are substantially longer than the dataset-wide average everywhere in the bundle. The average first-loop-prefix length also varies meaningfully by dataset, which indicates that not every loop is an immediate degeneration: some settings enter the repeated n-gram regime only after a long reasoning or coding prefix.

Rollout success under looping is dataset-dependent rather than uniform. Some tasks retain a noticeable fraction of correct answers even inside looped rollouts, which suggests the model can sometimes reach the right answer before it starts repeating. For \texttt{LiveCodeBench}, the native benchmark metrics should be read from the separate \texttt{pass@k} values rather than from the rollout-success column.

\end{document}
